% GENERATED by tools/coverage.py from the lesson plans -- do not edit by hand.
\subsection{Objectives per goal and week (generated)}
Each number is how many of that week's \enquote{The student can\ldots} objectives
serve the goal. Blank = not a focus that week (it may still be practiced).

{\small\setlength{\tabcolsep}{3pt}
\begin{longtable}{@{}lcccccccccccc@{}}
\toprule
Goal & 1 & 2 & 3 & 4 & 5 & 6 & 7 & 8 & 9 & 10 & B1 & B2\\ \midrule\endhead
J2.G1 & 2 & 1 &  &  & 1 &  &  & 1 & 1 &  &  & 1\\
J2.G2 &  & 2 & 1 &  & 2 &  &  & 1 & 1 &  & 1 & 1\\
J2.G3 &  &  & 2 & 2 & 2 &  &  & 2 & 1 &  &  & 1\\
J2.G4 &  &  &  &  &  & 2 & 2 &  &  & 1 &  & \\
J2.G5 & 1 & 1 & 1 & 1 & 1 & 1 & 1 & 1 & 1 & 1 & 1 & 1\\
J2.G6 &  & 1 &  &  & 1 &  &  & 1 & 2 & 1 & 1 & 2\\
J2.G7 & 1 & 1 & 1 & 1 &  & 1 & 1 & 1 & 1 & 1 & 1 & 1\\
\bottomrule\end{longtable}}

\subsection{All lesson objectives (generated)}
\begin{longtable}{@{}p{1.5cm}p{0.72\linewidth}p{1.6cm}@{}}
\toprule Week & Objective & Goals\\ \midrule\endhead
\multicolumn{3}{@{}l}{\textbf{1: Forever and Bounce}}\\
1.1 & The student can use a forever loop with \enquote{move 10 steps} and \enquote{if on edge, bounce} to keep a sprite moving until the red stop sign is clicked. & J2.G1\\
1.2 & The student can say how repeat and forever differ: a repeat stops after its number, a forever loop keeps going until someone presses stop. & J2.G1\\
1.3 & The student can find blocks placed under a forever loop, explain why they never run, and move them above the loop. & J2.G5\\
1.4 & The student can stop playing when the timer rings and say their own stop plan. & J2.G7\\
\midrule
\multicolumn{3}{@{}l}{\textbf{2: If Touching\ldots}}\\
2.1 & The student can put an \enquote{if touching \ldots{} then} block inside a forever loop so a sprite reacts every time it is touched. & J2.G1, J2.G2\\
2.2 & The student can build a catch game in which an apple falls from a hidden Star marker, and goes back to the Star when it touches the bowl or the edge. & J2.G2\\
2.3 & The student can find an if that sits outside the forever loop, explain that it checks only once, and move it inside. & J2.G5\\
2.4 & The student can give kind feedback on a partner's game with two stars and a wish. & J2.G6, J2.G7\\
\midrule
\multicolumn{3}{@{}l}{\textbf{3: Keep Score}}\\
3.1 & The student can make a variable called score \enquote{For all sprites} and show it on the stage. & J2.G3\\
3.2 & The student can set the score to 0 when the green flag is clicked and change it by 1 each time the bowl catches an apple. & J2.G2, J2.G3\\
3.3 & The student can find and fix a score that does not go back to 0 when the flag is clicked, and say how \enquote{set} and \enquote{change} differ. & J2.G5\\
3.4 & The student can say the scoring rules of their game before a friend plays, so that everyone knows them. & J2.G7\\
\midrule
\multicolumn{3}{@{}l}{\textbf{4: Surprise! Random}}\\
4.1 & The student can use \enquote{pick random} to roll a number in a variable and make a sprite hop that many times. & J2.G3\\
4.2 & The student can use \enquote{go to random position} and a score variable to make a game where a sprite pops up in surprise places. & J2.G3\\
4.3 & The student can find and fix a random number that is picked only once, before the forever loop, so it never changes. & J2.G5\\
4.4 & The student can say how to be a good sport when they win or lose a game. & J2.G7\\
\midrule
\multicolumn{3}{@{}l}{\textbf{5: Play Day and Show Me}}\\
5.1 & The student can be the computer: act out a script with an if-touching check and a score, and say the score at the end. & J2.G2, J2.G3\\
5.2 & The student can find and fix a bug in each of three broken games with \enquote{test, look, change one thing}. & J2.G5\\
5.3 & The student can build a small game with a forever loop, an if-touching check, a score that starts at 0 and goes up by 1, and a random surprise. & J2.G1, J2.G2, J2.G3\\
5.4 & The student can take turns as Driver and Helper, explain what their script does and say one thing they got better at. & J2.G6\\
\midrule
\multicolumn{3}{@{}l}{\textbf{6: Messages}}\\
6.1 & The student can make a start button that sends a message with \enquote{broadcast}, so that other sprites start with \enquote{when I receive}. & J2.G4\\
6.2 & The student can pass a message along a chain of three sprites so they take turns without counting seconds. & J2.G4\\
6.3 & The student can find and fix a message that nobody listens for and a \enquote{when I receive} with the wrong message name. & J2.G5\\
6.4 & The student can give credit in a game for sprites and sounds they did not make. & J2.G7\\
\midrule
\multicolumn{3}{@{}l}{\textbf{7: My Own Blocks}}\\
7.1 & The student can make a My Block with a name, such as \enquote{celebrate}, and use it in two different scripts. & J2.G4\\
7.2 & The student can split a job into named steps on recipe cards and say which group of steps each name stands for. & J2.G4\\
7.3 & The student can find and fix a My Block that was made but never used, by testing and looking for the block in the scripts. & J2.G5\\
7.4 & The student can make a game easier for everyone to play, with big text and things that differ by shape as well as color. & J2.G7\\
\midrule
\multicolumn{3}{@{}l}{\textbf{8: Start, Win and Lose}}\\
8.1 & The student can give their catch game a start that sets the score and misses to 0 and tells the player the goal, and a forever loop that keeps the game playing. & J2.G1, J2.G3\\
8.2 & The student can end the game with \enquote{if score = 10 then} You win! and \enquote{if misses = 3 then} Game over!, each followed by \enquote{stop all}. & J2.G2, J2.G3\\
8.3 & The student can find and fix a game that never ends or ends right at the start, by testing both endings. & J2.G5\\
8.4 & The student can say why good games end and do not trick players into \enquote{just one more} go, and write fair start, win and lose rules for their own arcade game. & J2.G6, J2.G7\\
\midrule
\multicolumn{3}{@{}l}{\textbf{9: Project Studio}}\\
9.1 & The student can plan their arcade game on the game plan card with a partner: the player, what to catch or avoid, how to win and how to lose. & J2.G6\\
9.2 & The student can build the core of their game one must-have at a time, testing each: a forever loop, an if-touching check, a score or misses variable, random and an ending. & J2.G1, J2.G2, J2.G3\\
9.3 & The student can watch a Player playtest their game, write a playtest bug list and fix one bug the Player found. & J2.G5, J2.G6\\
9.4 & The student can ask before sharing or showing someone else's game, and explain why our games stay with our family. & J2.G7\\
\midrule
\multicolumn{3}{@{}l}{\textbf{10: Arcade Showcase}}\\
10.1 & The student can present their arcade game in a one-minute demo and explain one design choice and why they made it, such as a change that came from a Player's wish. & J2.G6\\
10.2 & The student can run a last-minute bug check (green flag twice, play to win, play to lose) and fix one bug calmly. & J2.G5\\
10.3 & The student can show a message being sent and received, and a My Block being made and used, in their game or in a quick demo. & J2.G4\\
10.4 & The student can give kind sticker feedback in the arcade gallery, give credit for sprites and ideas they did not make, and keep the showcase for family only with nicknames on screen. & J2.G7\\
\midrule
\multicolumn{3}{@{}l}{\textbf{B1: Secret Codes Day}}\\
B1.1 & The student can plan a coded treasure hunt with a partner and change a clue after a player says it was too hard. & J2.G6\\
B1.2 & The student can code and decode a word with a letter-shift wheel, such as \enquote{shift each letter 3 forward}, and find a wrongly coded letter by checking one letter at a time. & J2.G5\\
B1.3 & The student can solve a logic puzzle with if-then clues and say which clue gave each answer. & J2.G2\\
B1.4 & The student can tell a fun secret code from a worried secret that they should tell a grown-up, and say that passwords stay private. & J2.G7\\
\midrule
\multicolumn{3}{@{}l}{\textbf{B2: Mini Game Jam}}\\
B2.1 & The student can come up with a game idea in 90 seconds and say it in one sentence: who you are, what you do, how you win. & J2.G6\\
B2.2 & The student can build a small game in one session with a forever loop, an if-then check and a score. & J2.G1, J2.G2, J2.G3\\
B2.3 & The student can let a partner play their jam game, listen to two stars and a wish, and change one thing. & J2.G5, J2.G6\\
B2.4 & The student can tell players how to win at the start and give credit for an idea they borrowed. & J2.G7\\
\midrule
\bottomrule\end{longtable}
